\chapter{平面中的向量与直线}\label{ch:g11:vect}

\omterm{def:g11:vect:vector}{向量}刻画位移，直线是其轨迹。本章的关键新工具是
\emph{\omterm{def:g11:vect:det}{行列式}}，它是一个数，
可以判断两个\omterm{def:g11:vect:vector}{向量}是否\omterm{def:g10:vectors:collinear}{共线}——由此出发，得出直线\omterm{def:g10:algebra:equation}{方程}，
并通过纯粹的计算解决三点\omterm{def:g10:vectors:collinear}{共线}与平行问题。
相同的思想将在\cref{ch:g12:space}中推广到空间。

\section{向量：回顾}

\begin{definition}[向量，坐标]\label{def:g11:vect:vector}
\emph{向量}\index{向量} $\vec u$ 表示一个位移：所有
方向、指向与长度均相同的箭头都表示同一个
向量。在\omterm{def:g10:coordgeom:system}{坐标系}中，$\vec u\,(x, y)$ 记录位移
（$x$ 为水平方向，$y$ 为竖直方向）；对于点 $A(x_A, y_A)$ 与
$B(x_B, y_B)$，
\[
\vect{AB}\,\bigl(x_B - x_A,\ y_B - y_A\bigr).
\]
向量按坐标相加，且 $\lambda\vec u$ 的坐标为
$(\lambda x, \lambda y)$。
\end{definition}

\begin{proposition}[中点]\label{prop:g11:vect:midpoint}
点 $M$ 为线段 $[AB]$ 的\omterm{prop:g10:coordgeom:midpoint}{中点}，其坐标为
$\left(\frac{x_A + x_B}{2},\ \frac{y_A + y_B}{2}\right)$。
\end{proposition}

\begin{proof}
$M$ 由 $\vect{AM} = \frac12\vect{AB}$ 刻画。按坐标分量，
$x_M - x_A = \frac12(x_B - x_A)$，于是 $x_M = \frac{x_A + x_B}{2}$，
对 $y_M$ 同理。
\end{proof}

\section{共线与行列式}

\begin{definition}[共线向量]\label{def:g11:vect:collinear}
若两个\omterm{def:g11:vect:vector}{向量}中的一个是另一个的数乘，则称它们为\emph{共线向量}\index{共线向量}，
即 $\vec v = \lambda \vec u$，其中
$\lambda \in \R$（或 $\vec u = \vec 0$）。几何上看，它们方向相同，
或其中之一为零\omterm{def:g11:vect:vector}{向量}。
\end{definition}

\begin{definition}[行列式]\label{def:g11:vect:det}
\omterm{def:g11:vect:vector}{向量} $\vec u\,(x, y)$ 与 $\vec v\,(x', y')$ 的
\emph{行列式}\index{行列式}为
\[
\det(\vec u, \vec v) =
\begin{vmatrix} x & x' \\ y & y' \end{vmatrix}
= x y' - x' y .
\]
\end{definition}

\begin{theorem}[共线判定准则]\label{thm:g11:vect:collinearity}
两个\omterm{def:g11:vect:vector}{向量} $\vec u$ 与 $\vec v$ \omterm{def:g10:vectors:collinear}{共线}，当且仅当
$\det(\vec u, \vec v) = 0$。
\end{theorem}

\begin{proof}
（$\Rightarrow$）若 $\vec v = \lambda\vec u$，则 $x' = \lambda x$，且
$y' = \lambda y$，于是
$\det(\vec u, \vec v) = x(\lambda y) - (\lambda x)y = 0$；若
$\vec u = \vec 0$，则\omterm{def:g11:vect:det}{行列式}显然为 $0$。

（$\Leftarrow$）设 $xy' - x'y = 0$，且 $\vec u \neq \vec 0$，不妨设
$x \neq 0$（$y \neq 0$ 的情形同理）。令
$\lambda = \frac{x'}{x}$，则由构造 $x' = \lambda x$，又
关系式 $xy' = x'y = \lambda x y$ 两边除以 $x \neq 0$，得
$y' = \lambda y$，因此 $\vec v = \lambda \vec u$。
\end{proof}

\begin{example}\label{ex:g11:vect:collinear}
$\vec u\,(3, -2)$ 与 $\vec v\,(-6, 4)$：
$\det = 3 \times 4 - (-6) \times (-2) = 12 - 12 = 0$，故\omterm{def:g10:vectors:collinear}{共线}（事实上
$\vec v = -2\vec u$）。$\vec u\,(3, -2)$ 与 $\vec w\,(2, 1)$：
$\det = 3 + 4 = 7 \neq 0$，故不\omterm{def:g10:vectors:collinear}{共线}。
\end{example}

\begin{method}[三点共线]\label{met:g11:vect:alignment}
$A$、$B$、$C$ 三点\omterm{def:g10:vectors:collinear}{共线}，当且仅当 $\vect{AB}$ 与 $\vect{AC}$
\omterm{def:g10:vectors:collinear}{共线}：计算这两个\omterm{def:g11:vect:vector}{向量}，并检验
$\det\bigl(\vect{AB}, \vect{AC}\bigr) = 0$。
\end{method}

\section{直线的一般式方程}

\begin{definition}[方向向量]\label{def:g11:vect:direction}
称直线 $d$ 的\emph{方向向量}\index{方向向量}为满足下述条件的任一
非零\omterm{def:g11:vect:vector}{向量} $\vec u$，使得 $\vect{AB}$ 与 $\vec u$ \omterm{def:g10:vectors:collinear}{共线}，
其中 $A, B$ 为 $d$ 上的任意两点。
\end{definition}

\begin{theorem}[一般式方程]\label{thm:g11:vect:cartesian}
\omterm{def:g11:vect:direction}{方向向量}为 $\vec u\,(-b, a)$ 的直线，
可写成如下形式的\omterm{def:g10:algebra:equation}{方程}：
\[
ax + by + c = 0
\qquad (a, b) \neq (0, 0).
\]
反之，每个这样的\omterm{def:g10:algebra:equation}{方程}都表示一条\omterm{def:g11:vect:direction}{方向向量}为
$\vec u\,(-b, a)$ 的直线。
\end{theorem}

\begin{proof}
设直线 $d$ 过点 $A(x_A, y_A)$，\omterm{def:g11:vect:direction}{方向向量}为 $\vec u\,(-b, a)$。
点 $M(x, y)$ 在 $d$ 上，当且仅当 $\vect{AM}$ 与
$\vec u$ \omterm{def:g10:vectors:collinear}{共线}，即由\cref{thm:g11:vect:collinearity}可得
\[
0 = \det\bigl(\vect{AM}, \vec u\bigr)
= (x - x_A)\,a - (-b)(y - y_A)
= ax + by - (a x_A + b y_A),
\]
这是所述形式的\omterm{def:g10:algebra:equation}{方程}，其中 $c = -(a x_A + b y_A)$。
反之，对\omterm{def:g10:algebra:equation}{方程} $ax + by + c = 0$，不妨设 $b \neq 0$，则其解为
$\left(x, -\frac{a x + c}{b}\right)$。将其中两点相减，
可知任意两解的连线都与 $(-b, a)$ \omterm{def:g10:vectors:collinear}{共线}，并且
总存在一个解，解的全体即为过该解、以 $(-b, a)$ 为方向的直线。
\end{proof}

\begin{method}[过两点的直线]\label{met:g11:vect:twopoints}
求直线 $(AB)$ 的\omterm{def:g10:algebra:equation}{方程}：计算\omterm{def:g11:vect:direction}{方向向量}
$\vect{AB}\,(x_B - x_A, y_B - y_A)$；使其与 $(-b, a)$ 对应，即取
$a = y_B - y_A$，$b = -(x_B - x_A)$；再代入以确定 $c$，
即把 $A$ 的坐标代入。
\end{method}

\begin{example}\label{ex:g11:vect:twopoints}
过 $A(1, 2)$ 和 $B(4, 3)$ 的直线：$\vect{AB}\,(3, 1)$，故
$a = 1$，$b = -3$，\omterm{def:g10:algebra:equation}{方程}为 $x - 3y + c = 0$。代入 $A$：
$1 - 6 + c = 0$，故 $c = 5$。\omterm{def:g10:algebra:equation}{方程}为 $x - 3y + 5 = 0$。用 $B$ 检验：
$4 - 9 + 5 = 0$.
\end{example}

\begin{remark}
当 $b \neq 0$ 时，\omterm{def:g10:algebra:equation}{方程}可解出 $y$，得 $y = mx + p$，
\omterm{def:g10:reffunc:affine}{斜率}为 $m = -\frac ab$。直线的一般式 $ax + by + c = 0$
适用范围更广，它还包括竖直线（$b = 0$），这类直线没有\omterm{def:g10:reffunc:affine}{斜率}。
\end{remark}

\begin{omfigure}
\begin{tikzpicture}
\begin{axis}[omaxis,
  width=8.6cm, height=5.8cm,
  xmin=-1.6, xmax=5.6, ymin=-0.9, ymax=4.4,
  xtick={1,4}, ytick={2,3}]
  \addplot[omDef, thick, domain=-1.4:5.4] {(x+5)/3};
  \fill (axis cs:1,2) circle (1.5pt)
    node[above left, font=\small] {$A$};
  \fill (axis cs:4,3) circle (1.5pt)
    node[above left, font=\small] {$B$};
  \draw[-Stealth, omThm, thick] (axis cs:1,2) -- (axis cs:2.5,2.5)
    node[midway, below right, font=\small] {$\vec u$};
\end{axis}
\end{tikzpicture}

{\small \cref{ex:g11:vect:twopoints}中的直线 $x - 3y + 5 = 0$，其
\omterm{def:g11:vect:direction}{方向向量} $\vec u = \vect{AB}$（红色，此处按 $\frac12$ 缩放），
以及建立\omterm{def:g10:algebra:equation}{方程}时所用的两点。}
\end{omfigure}

\section{直线的参数表示}

\begin{definition}[直线的参数表示]\label{def:g11:vect:parametric}
过 $A(x_A, y_A)$、以 $\vec u\,(\alpha, \beta)$ 为\omterm{def:g11:vect:direction}{方向向量}的直线
是如下点的集合：
\[
M\bigl(x_A + t\alpha,\ y_A + t\beta\bigr), \qquad t \in \R .
\]
数 $t$ 称为\emph{参数}：$t$ 的每个取值对应直线上的一个
点，好像以恒定速度 $\vec u$ 沿直线运动一样。
\end{definition}

\begin{example}
\cref{ex:g11:vect:twopoints}中的直线也可表示为
$\{(1 + 3t,\ 2 + t) : t \in \R\}$：当 $t = 0$ 时得 $A$，$t = 1$ 时得 $B$。
消去 $t$（由 $t = y - 2$，故 $x = 1 + 3(y-2)$），可得
$x - 3y + 5 = 0$.
\end{example}

\section{两直线的位置关系}

\begin{proposition}[平行或相交]\label{prop:g11:vect:positions}
\omterm{def:g11:vect:direction}{方向向量}分别为 $\vec u$ 和 $\vec v$ 的两直线平行，当且仅当
$\det(\vec u, \vec v) = 0$。若用\omterm{def:g10:algebra:equation}{方程}形式表示，则 $d : ax + by + c =
0$ 与 $d' : a'x + b'y + c' = 0$ 平行，当且仅当
$ab' - a'b = 0$。不平行的两直线恰交于一点。
\end{proposition}

\begin{proof}
平行即\omterm{def:g11:vect:direction}{方向向量}\omterm{def:g10:vectors:collinear}{共线}，
由\cref{thm:g11:vect:collinearity}可知。取 $\vec u\,(-b, a)$ 与
$\vec v\,(-b', a')$，则\omterm{def:g11:vect:det}{行列式}为 $(-b)a' - (-b')a = ab' - a'b$。
若两直线不平行，则联立两个\omterm{def:g10:algebra:equation}{方程}
（用代入法或加减消元法）解得唯一解：此时\omterm{def:g10:algebra:equation}{方程}组
类似于\omterm{def:g11:vect:det}{行列式}不为零的 $2 \times 2$ \omterm{def:g10:algebra:equation}{方程}组。
\end{proof}

\begin{method}[两直线的交点]\label{met:g11:vect:intersection}
解由这两个\omterm{def:g10:algebra:equation}{方程}组成的\omterm{def:g10:algebra:equation}{方程}组：在一个\omterm{def:g10:algebra:equation}{方程}中解出一个未知数，
并代入另一个\omterm{def:g10:algebra:equation}{方程}，或将两\omterm{def:g10:algebra:equation}{方程}加减以消去一个未知数。
所得的点都须代入\emph{两个}\omterm{def:g10:algebra:equation}{方程}进行检验。
\end{method}

\begin{example}\label{ex:g11:vect:intersection}
求 $d: x - 3y + 5 = 0$ 与 $d': 2x + y - 4 = 0$ 的交点。由 $d'$ 得：
$y = 4 - 2x$。代入 $d$：
\[
x - 3(4 - 2x) + 5 = 0
\iff x - 12 + 6x + 5 = 0
\iff 7x = 7
\iff x = 1,
\]
于是 $y = 2$。两直线交于 $(1, 2)$。在 $d$ 中检验：$1 - 6 + 5 = 0$。
\end{example}

\section{练习}

\begin{exercise}[$\star$]\label{exo:g11:vect:1}
已知 $A(2, -1)$、$B(5, 3)$ 与 $C(-1, 1)$，求
$\vect{AB}$、$\vect{AC}$、$\vect{AB} + \vect{AC}$ 与 $2\vect{AB} -
3\vect{AC}$ 的坐标，以及线段 $[BC]$ \omterm{prop:g10:coordgeom:midpoint}{中点}的坐标。
\end{exercise}

\begin{exercise}[$\star$]\label{exo:g11:vect:2}
判断下列各组\omterm{def:g11:vect:vector}{向量}是否\omterm{def:g10:vectors:collinear}{共线}。
\[
\vec u\,(4, -6) \text{ 与 } \vec v\,(-2, 3); \qquad
\vec u\,(2, 5) \text{ 与 } \vec v\,(3, 7); \qquad
\vec u\,(0, 3) \text{ 与 } \vec v\,(0, -8).
\]
\end{exercise}

\begin{exercise}[$\star$]\label{exo:g11:vect:3}
求过点 $A(2, 1)$、\omterm{def:g11:vect:direction}{方向向量}为
$\vec u\,(3, -1)$ 的直线的普通\omterm{def:g10:algebra:equation}{方程}，再求过 $B(0, 4)$ 和
$C(2, 0)$ 的直线的普通\omterm{def:g10:algebra:equation}{方程}。
\end{exercise}

\begin{exercise}[$\star$]\label{exo:g11:vect:4}
对直线 $d: 3x - 2y + 6 = 0$，求其一个\omterm{def:g11:vect:direction}{方向向量}、
与两坐标轴的交点，
并判断点 $P(2, 6)$ 与 $Q(1, 4)$ 是否在直线 $d$ 上。
\end{exercise}

\begin{exercise}[$\star\star$]\label{exo:g11:vect:5}
点 $A(1, 2)$、$B(4, 8)$、$C(-2, -4)$ 是否\omterm{def:g10:vectors:collinear}{共线}？
再判断点 $D(0, 1)$、$E(2, 4)$、$F(5, 9)$ 是否\omterm{def:g10:vectors:collinear}{共线}。
\end{exercise}

\begin{exercise}[$\star\star$]\label{exo:g11:vect:6}
求交点（若存在）：
\[
d: 2x + y - 7 = 0 \ \text{ 与 }\ d': x - y + 1 = 0;
\qquad
e: x - 2y + 3 = 0 \ \text{ 与 }\ e': -2x + 4y + 1 = 0 .
\]
\end{exercise}

\begin{exercise}[$\star\star$]\label{exo:g11:vect:7}
写出直线 $d: 2x - y + 3 = 0$ 的参数\omterm{def:g10:algebra:equation}{方程}，
并求下列直线的普通\omterm{def:g10:algebra:equation}{方程}：
$\bigl\{(1 + 2t,\ -3 + t) : t \in \R\bigr\}$。
\end{exercise}

\begin{exercise}[$\star\star$]\label{exo:g11:vect:8}
求过点 $P(1, -2)$ 且与
$d: 4x - 3y + 2 = 0$ 平行的直线\omterm{def:g10:algebra:equation}{方程}，以及 $m$ 为何值时，直线
$mx + 2y - 5 = 0$ 与 $d$ 平行。
\end{exercise}

\begin{exercise}[$\star\star$]\label{exo:g11:vect:9}
设 $A(-1, 0)$、$B(3, 2)$、$C(1, -4)$，求三角形 $ABC$ 中
自顶点 $C$ 引出的\emph{中线}（即连接 $C$
与 $[AB]$ \omterm{prop:g10:coordgeom:midpoint}{中点}的直线）的普通\omterm{def:g10:algebra:equation}{方程}。
\end{exercise}

\begin{exercise}[$\star\star$]\label{exo:g11:vect:10}
参数 $m$ 取何值时，\omterm{def:g11:vect:vector}{向量}
$\vec u\,(m, 2)$ 与 $\vec v\,(3, m - 1)$ \omterm{def:g10:vectors:collinear}{共线}？
\end{exercise}

\begin{exercise}[$\star\star\star$]\label{exo:g11:vect:11}
设 $ABCD$ 为平行四边形，$I$ 为 $[AB]$ 的\omterm{prop:g10:coordgeom:midpoint}{中点}，$J$ 为
满足 $\vect{DJ} = \frac23 \vect{DI}$ 的点。取\omterm{def:g10:coordgeom:system}{坐标系}使
$A(0,0)$、$B(1,0)$、$D(0,1)$（从而 $C(1,1)$）：
\begin{enumerate}
  \item 求 $I$ 与 $J$ 的坐标。
  \item 证明 $A$、$J$、$C$ 三点\omterm{def:g10:vectors:collinear}{共线}。（一个有趣的事实：直线
        $(DI)$ 与对角线 $(AC)$ 的交点恰在对角线的三分之一处。）
\end{enumerate}
\end{exercise}

\section{问题：碰撞航线与巧选坐标}

\begin{problem}[{周末问题——路径相交不等于相撞：
海上的最近接近，以及通过
选取恰当坐标系证明的定理}]
\label{pb:g11:vect:1}
两艘船的直线航线相交——两船一定会相撞吗？并非如此，
除非两船\emph{同时}到达交点：路径是点的集合，
运动则是参数化的，混淆二者曾使真实的船只沉没。
本题将搭建一个小型船舶交通服务系统，
基于参数直线（见\cref{def:g11:vect:parametric}），然后转向纯几何问题，
并利用坐标的另一大优势：\emph{巧妙选取}坐标，
如\cref{exo:g11:vect:11}所示，使经典定理
“自动”算出结果。

\medskip
\noindent\textbf{第一部分——直线的三种表示法。}
\begin{enumerate}
  \item 写出过点
        $A(1, 2)$、\omterm{def:g11:vect:direction}{方向向量}为 $\vec u(3, -1)$ 的直线的参数\omterm{def:g10:algebra:equation}{方程}，然后
        消去参数，得到其普通\omterm{def:g10:algebra:equation}{方程}。
  \item 对直线 $2x - 5y + 3 = 0$，读出一个\omterm{def:g11:vect:direction}{方向向量}
        （见\cref{thm:g11:vect:cartesian}），求出直线上一点，
        并写出其参数\omterm{def:g10:algebra:equation}{方程}。
  \item 求第 1 题与第 2 题中两直线的交点
        （见\cref{met:g11:vect:intersection}）。
  \item 利用\omterm{def:g11:vect:det}{行列式}
        （见\cref{thm:g11:vect:collinearity}）验证那两条直线必定相交；
        再证明 $2x - 5y + 3 = 0$ 与
        $-4x + 10y + 1 = 0$ 平行且不重合。
  \item 点 $A(2, 1)$、$B(5, 3)$、$C(11, 7)$ 是否\omterm{def:g10:vectors:collinear}{共线}
        （见\cref{met:g11:vect:alignment}）？
\end{enumerate}

\medskip
\noindent\textbf{第二部分——船舶交通服务。}
在时刻 $t$（单位：小时）时，船 $P$ 位于 $(2t,\ 3 + t)$，船 $Q$
位于 $(10 - t,\ 2t - 1)$（距离单位为海里）。
\begin{enumerate}[resume]
  \item 求两条\emph{路径}的普通\omterm{def:g10:algebra:equation}{方程}，
        并求两条路径相交的位置。
  \item 两船各在何时经过交点？两船是否
        相撞？用一句话给出一般性的提醒：
        路径相交，不等于运动相遇。
  \item 求两船在任一时刻的距离平方 $D^2(t)$，
        并化为关于 $t$ 的\omterm{def:g11:quad:quadratic}{二次函数}。问：
        两船在 $t$ 为何值时距离最近？最近距离为多少
        （用\cref{pb:g11:quad:1}中的顶点法）？
  \item 规定两船间距至少为 $1$ 海里。
        判断该规定是否被违反？若是，则解不等式 $D^2(t) < 1$，
        并求违反规定的时长。
  \item 解释为何对\emph{任意}两艘沿直线航行、速度恒定的船，
        $D^2(t)$ 总是
        $t$ 的\omterm{def:g11:quad:quadratic}{二次函数}——因此“最近接近点”
        总是一个顶点计算问题。
  \item 拦截问题：一艘巡逻艇于 $t = 0$ 时从原点出发，
        以恒定速度 $(a, b)$ 运动，并须与船 $P$
        在 $t = 3$ 时恰好相遇。求所需的
        $(a, b)$ 及巡逻艇的速率。
\end{enumerate}

\medskip
\noindent\textbf{第三部分——巧妙的坐标。}
采用\cref{exo:g11:vect:11}的框架：平行四边形
$ABCD$ 变为 $A(0,0)$、$B(1,0)$、$C(1,1)$、$D(0,1)$；设 $I$
为 $[AB]$ 的\omterm{prop:g10:coordgeom:midpoint}{中点}，$K$ 为 $[CD]$ 的\omterm{prop:g10:coordgeom:midpoint}{中点}。
\begin{enumerate}[resume]
  \item 求直线 $(DI)$ 的直角坐标\omterm{def:g10:algebra:equation}{方程}。
  \item 求 $(DI)$ 与对角线 $(AC)$ 的交点，并验证
        \cref{exo:g11:vect:11} 中的结论：该交点恰好
        位于对角线的三分之一处。
  \item 再看直线 $(BK)$：求其\omterm{def:g10:algebra:equation}{方程}，并与
        $(AC)$ 求交点，从而得到完整的经典定理：
        \emph{塞瓦线 $(DI)$ 与 $(BK)$ 三等分对角线
        $(AC)$}。
  \item 为何可以在这一特殊坐标框架中证明该定理？
        （坐标的选取保持了什么，命题涉及什么？）
        请用
        两到三句话作答。
  \item 再练习一个三等分问题：设 $E$ 为
        $[BC]$ 的\omterm{prop:g10:coordgeom:midpoint}{中点}。直线 $(AE)$
        与\emph{另一条}对角线 $(DB)$ 相交于何处？
\end{enumerate}

\medskip
\noindent\textbf{第四部分——位置谜题。}
\begin{enumerate}[resume]
  \item 判断三条直线 $x + y = 4$、\ $2x - y = 2$、\
        $x - 2y = -2$ 是否共点？（先求其中两条的交点，
        再检验第三条。）
  \item 对每个\omterm{def:g10:numbers:sets}{实数} $m$，设 $d_m$ 为直线
        $mx - y + 2 - 3m = 0$。证明\emph{每一条} $d_m$
        都经过一个定点。（将关于
        $m$ 的项归并。）
  \item 在族 $(d_m)$ 中，哪一条与
        $y = 2x + 7$ 平行？哪一条经过原点？
  \item 总结——直线的三种形式：直角坐标\omterm{def:g10:algebra:equation}{方程}
        （代入一次即可判断点是否在直线上）、斜截式（\omterm{def:g10:reffunc:affine}{斜率}与\omterm{def:g10:functions:graph}{图像}一目了然）、
        参数\omterm{def:g10:algebra:equation}{方程}（描述运动、计时与拦截）；
        \omterm{def:g11:vect:det}{行列式}作为判断平行的工具；
        以及选取恰当的坐标框架这一定理的“生产工厂”。
        各用一句话作答。
\end{enumerate}
\end{problem}
